\section{Open questions for follow-up}

The following five questions emerged from a careful re-read of the paper and remain \textbf{genuinely open} after this replication effort. None was resolvable from the manuscript alone. Each is stated with (i) a basis grounded in a specific passage or missing artifact, and (ii) a concrete, executable next step.

\subsection*{Q1 --- What numerical value of $D_f$ (the ROS diffusion/recombination constant) is implied by the ReaxFF trajectories?}
\textbf{Basis:} The rate equations (Eqs~1--2) contain $D_f$ as the sole non-source coupling constant, and Figs~6--7 are shown in dimensionless units without a physical time axis. The paper implies $D_f$ could be derived from the MD but never gives a number. Without $D_f$ the shapes in Figs~6--7 cannot be placed on a physical time axis (ns? µs? ms?), which is what determines whether NROS chains actually survive the inter-pulse gap at 40~Gy/s FLASH clinical rates.\\
\textbf{Next steps:} Ask the corresponding author for the LAMMPS ReaxFF trajectory files or a mean-square-displacement time-series for $\cdot$OH; alternatively, run a small LAMMPS-ReaxFF MSD calculation on liquid water + $\cdot$OH at 300~K using the Rahaman 2011 FF and estimate $D_f = D_{\text{OH}}/(6\langle r_{\text{react}}^2\rangle)$ with $r_{\text{react}} \approx 3$~\AA. Cross-check against Elliot \& Bartels 2009 water-radiolysis reaction-diffusion tables.

\subsection*{Q2 --- Does $\mu_1(N_{O_2}, N_{H_2O})$ exhibit the non-monotonic O\textsubscript{2}-dependence claimed in \S III?}
\textbf{Basis:} The paper asserts (text-only, no figure) an oxygen optimum for FLASH sparing at $\sim$4--5\% O\textsubscript{2} that reverses above $\sim$10--15\%. This is a strong biological claim (it predicts hyperoxic tissue would be \emph{harmed} by FLASH), but no $\Delta$O\textsubscript{2} vs $[\text{O}_2]_0$ curve is provided --- only qualitative MD snapshots (Figs~2--5). This is a specific, falsifiable, clinically relevant prediction that the paper leaves quantitatively empty.\\
\textbf{Next steps:} Re-read \S III of the paper to identify whether $\mu_1$ is defined implicitly or explicitly; if implicit, request the source-function code from authors. Then sweep the rate equations at fixed $G$, varying $\mu_1$'s [O\textsubscript{2}] dependence, to see whether the optimum re-emerges. Cross-check against the Zhu et al.\ 2020 \emph{Radiother Oncol} data on hyperoxia-plus-FLASH.

\subsection*{Q3 --- How sensitive is the $N_2(\text{FLASH})/N_2(\text{CDR})\approx 2\times$ ratio to the pulse-parameter choice?}
\textbf{Basis:} The authors admit the pulse parameters are chosen to \emph{illustrate the method}, not to match Gy/s. Our reproduction confirms the $\approx 2\times$ ratio at $t=100$~s in these units, but a factor-of-2 that depends entirely on an arbitrary parameter selection is not a robust prediction. If the ratio is highly sensitive to $(G_1\tau_1)/(G_2\tau_2)$ at fixed integral, or to the ratio $\tau_2/\tau_1$, the ``FLASH doubles NROS'' headline collapses into ``FLASH-vs-CDR NROS ratio is model-tunable''.\\
\textbf{Next steps:} Extend the reproduction script to sweep $(G_1, \tau_1, G_2, \tau_2)$ with fixed $\int G\, dt = 1$ and plot $N_2(\text{FLASH})/N_2(\text{CDR})$ as a heat-map. Identify the region where the ratio exceeds $\sim$1.5$\times$. Compare that region to real UHDR-vs-CDR clinical dose-rate ratios ($\sim 10^3$--$10^4$).

\subsection*{Q4 --- Do the $\alpha$-oxygen / $\cdot$OH--$\cdot$OH coalescence products (NROS) have an experimentally-detectable spectroscopic signature?}
\textbf{Basis:} The paper introduces ``NROS chains'' as a new class of transient radical clusters and proposes they are the FLASH-differentiating species. But it does not identify a UV-Vis, EPR, or Raman fingerprint for them, nor cite an experiment showing them. If NROS are only visible in MD but not in the wet lab, the mechanism is not falsifiable outside computation.\\
\textbf{Next steps:} Search literature for experimental detection of transient $\cdot$OH clusters or trihydrogen radicals in pulsed radiolysis (LaVerne, Buxton EPR studies). Contact the authors for the specific chemical identity (Lewis structures, expected lifetimes) of the top-3 most abundant NROS species produced in their ReaxFF trajectories. Propose an EPR spin-trapping experiment (DMPO or PBN) at UHDR vs CDR to look for a differential adduct signature.

\subsection*{Q5 --- Why is the Rahaman 2011 glycine ReaxFF appropriate for a DNA + water + O\textsubscript{2} + radical system?}
\textbf{Basis:} The paper cites only ref 17 for its ReaxFF force field, but that force field was trained on glycine + water, not on nucleobases, sugar-phosphate backbone, or high concentrations of dissolved O\textsubscript{2} and radical species. ReaxFF is notoriously non-transferable --- using a peptide FF for DNA + radical chemistry is a strong extrapolation. The MD ``chemistry'' observed in Figs 1--5 ($\cdot$OH coupling, $\alpha$-O intermediates, chain formation) is force-field-dependent, and no sensitivity test or benchmark against alternative FFs (Ashraf--van Duin ReaxFF, ab initio) is presented.\\
\textbf{Next steps:} Re-read the ReaxFF literature (van Duin 2001 onwards) to confirm transferability caveats. Ask the authors whether they benchmarked chain-formation events against a smaller CPMD/DFT calculation, and whether they cross-validated with the Monti/Ashraf DNA-water ReaxFF (2013). Ideally rerun the last stage with a second FF and compare $\cdot$OH$\to\cdot$OH coupling rates as a sensitivity test.
